% Evaluation: the executed campaigns. Every figure in this section is copied
% from a retained record, and the sentence or table note that carries it names
% that record. Paths are repository-relative for files in dispat-research and
% absolute otherwise. Campaigns are never pooled; failed attempts are reported.
% \evpath, \evhash and \evRQ are defined in macros.tex.

\section{Evaluation: the executed campaigns}
\label{sec:evaluation}

The proofs of \cref{sec:guarantees} say what the algorithm does; this section
reports what the reference implementation~\cite{dispat} was observed to do, in
the campaigns that were actually executed, with the binary, the environment and
the counting unit of each one stated separately. Five questions organise the
account. \evRQ{1} asks whether the consumer-relative admission window retains
the residual obligations that a source-relative window loses. \evRQ{2} asks
whether versions and completion records stay stable through repeated failure,
repair and identity conflict, and how far the plans agree with an independent
model. \evRQ{3} asks whether execution respects artefact availability and CI
gates. \evRQ{4} asks what planning costs as the graph, the intent and the
history grow. \evRQ{5}, new in this report, asks what the distributed execution
profile of \spec{28} measured on two real build graphs, and what it did not.

The discipline of the section is that every number is transcribed from a
retained file that the sentence or the table note names, that no campaign is
pooled with another, and that failed and discarded attempts appear beside the
completed ones. Deterministic designed matrices are reported as counts with
their denominators and receive no $p$-values.

\subsection{Evidence boundaries and units}
\label{sec:eval:boundaries}

Seven bodies of evidence are used. \Cref{tab:eval:boundaries} pins each one to
its binary, its environment and its counting unit. The units differ and are
not interchangeable: a \emph{history} is one generated repository planned once;
an \emph{episode} is one fixture driven through fault, retry and repair; a
\emph{state} or a \emph{transition} is an enumeration unit of the finite model;
a \emph{command record} is one invocation the verifier replayed; a
\emph{predicate} is one verifier check applied repeatedly to fixtures; a
\emph{release} is one \code{dispat release} from invocation to exit under
GNU \code{time}. Predicate counts are implementation checks, not independent
observations, and are never summed across campaigns.

\begin{center}
\footnotesize
\begin{longtable}{@{}L{2.4cm}L{4.2cm}L{3.4cm}L{4.3cm}@{}}
\caption{Evidence boundaries. Sources: \protect\evpath{experiments/README.md} and
\protect\evpath{experiments/results/compare-v140-v160.md} (historical differential);
\protect\evpath{output/experiments-v182-summary.json} and
\protect\evpath{experiments/results/v182-20260907/README.md} (1.8.2);
\protect\evpath{paper/saner/evidence/EXPERIMENTS-1.9.0-20260909.md} (1.9.0);
\protect\evpath{paper/saner/evidence/EXPERIMENTS-1.10.0-20260909.md} and
\protect\evpath{experiments/saner-2027/campaign-1.10.0-20260909/summary.json} (1.10.0);
\protect\evpath{experiments/saner-2027/campaign-mechanisms-20260910/RESULTS.md}
(mechanisms); \protect\evpath{/Users/yohimik/Projects/linux-build-test/dispat/gce/PARITY-20260922.md}
and \protect\evpath{PARITY-20260922-round2.md} beside it (kernel);
\protect\evpath{.ci/dispat/gce/REPORT.md} on branch \code{origin/dispat-gce} of
\protect\evpath{/Users/yohimik/Projects/llvm-project} (LLVM).}
\label{tab:eval:boundaries}\\
\toprule
Campaign & Binary under test & Environment & Counting unit and record \\
\midrule
\endfirsthead
\toprule
Campaign & Binary under test & Environment & Counting unit and record \\
\midrule
\endhead
Historical differential, from 30 August 2026 & Released \code{dispat 1.4.0
(linux\_arm64)} from the official \evpath{yohimik/dispat-alpine:1.4.0} image
(release commit \code{5e3a2337}); repeated on 1.6.0 (\code{f9c2acf2}), 1.7.0
and 1.7.2 & Docker on an Apple-silicon host; one container at a time &
History (generated instance); \evpath{experiments/results/v140}, \evpath{v160},
\evpath{v170}, \evpath{v172} and the \evpath{compare-*.md} notes beside them \\
1.8.2 E1 to E5, 7 September 2026 & \code{dispat 1.8.2 (linux\_arm64)}, CLI
commit \evhash{a7011cbcb9e88e52e7336f6fc8f1b207a8e7a311}, image
\evpath{yohimik/dispat-alpine@sha256:}\allowbreak\evhash{4c0ad4b8268dd10dbb136008798c1b07ffc1c0c7c534b8465a9f5edcab0e7bfd}
& Docker, two CPUs and 4~GiB per container, contended host & History (E1),
repetition (E2), release row (E3), episode and timing cell (E4), check (E5);
\evpath{experiments/results/v182-20260907}, 753 files under its manifest \\
1.9.0 companion, 9 September 2026 & \code{dispat 1.9.0 (linux\_arm64)}, SHA-256
\evhash{be3370c472cdbfe66fe4256d3fc1ff78ff19bdd5f041ef222187a52767fa4f35},
runner image \evpath{sha256:}\allowbreak\evhash{97363c9dbac292599987e66241d6da56abd7540df965b24bc20d20573fcac31f}
& Docker Desktop Linux ARM64, two CPUs and 4~GiB, no network except private
registries & History, episode, execution, CI repetition, predicate;
\evpath{experiments/saner-2027/campaign-next-release-prepared/runs}, one failed
and one completed attempt \\
1.10.0 companion, 9 September 2026 & \code{dispat 1.10.0 (linux\_arm64)},
SHA-256 \evhash{03d1ba5fbaef35b69582e9489f40b5438f81ae6926e02ea6d02716b5b73d92fd},
runner image \evpath{sha256:}\allowbreak\evhash{6f92932fbba5d33d61d777073f5fc155133063f4c3d4b2f950e3e5b9ed9a263e}
& As above & As above; \evpath{experiments/saner-2027/campaign-1.10.0-20260909},
archive of 9,157 files \\
Mechanism matrix, 9 September 2026 & The same 1.10.0 binary and runner image
as the row above & Docker, network disabled, two CPUs, 4~GiB & Episode,
state, transition; \evpath{experiments/saner-2027/campaign-mechanisms-20260910},
16,464 fixture files archived \\
Kernel farm rounds, 22 September 2026 & Round 1: a development build, SHA-256
\evhash{a7172a7feb21e27ce882050120b9e04df1e6d4b4ee5f3341ab8a99ba155ba70a}.
Round 2: the feature branch at \code{f0c9e959}, static linux/amd64, SHA-256
\evhash{63441cc66f68d07b5efbf9a42edd5559304d153d48b82b1dfb037349e20606ed} &
Google Compute Engine~\cite{gce}: \code{c3-standard-22} and
\code{t2d-standard-8}, region us-central1 & Release (GNU \code{time} elapsed),
per-cell script seconds; \evpath{dispat/gce} of the fork at \code{c699bf3d}, copied under \evpath{experiments/distributed-20260922/evidence/kernel-gce} \\
LLVM pair, 22 September 2026 & Attempts 1 to 8: the feature branch at
\code{f0c9e959}; attempt 9: \code{1.11.0-rc.4+main.272772cb} as printed in its
release log & Google Compute Engine: \code{e2-standard-4} orchestrator, two
\code{n2-custom-14-49152} workers, zone us-central1-a & Release; the fork's
\evpath{.ci/dispat/gce/REPORT.md} and the retained logs under
\evpath{experiments/distributed-20260922/evidence/llvm-gce/logs} \\
\bottomrule
\end{longtable}
\end{center}

None of these campaigns measures a release-candidate feature outside
\evRQ{5}. The parser in every campaign through 1.10.0 is CCME 2.0.0
(\evpath{paper/saner/evidence/EXPERIMENTS-1.10.0-20260909.md}); the kernel and
LLVM rounds run the \spec{28} execution profile of the 1.11.0 release candidate
and are the only evidence about it.

\subsection{RQ1: the consumer-window mechanism}
\label{sec:eval:rq1}

The consumer-relative window of \cref{def:window} is the mechanism behind
\G{2}. The mechanism campaign tests it two ways, from
\evpath{experiments/saner-2027/campaign-mechanisms-20260910/PROTOCOL.md}: a
finite projection of the model, and a counterfactual computed on the
completion states the binary actually produced.

\paragraph{The finite projection.}
The projection enumerates every DAG consistent with one fixed order of four
named vertices, each of four single intent sources, every completion set
reachable by repeatedly recording an enabled package, and every subset of the
remaining packages designated as failed. It checks exact residual obligations,
no widening, no double ledger delivery, and strict decrease exactly when
progress occurs. It then applies the source-window mutant, which discards a
unit when its originating package completes, to the same states. The counts,
from the \code{executor.model} object of
\evpath{experiments/saner-2027/campaign-mechanisms-20260910/summary.json}, are
64 graphs, 256 source instances, 772 reachable completion states, 2,270
failure-subset transitions of which 1,054 made progress, 0 projection
violations, and 260 states in which the source-window mutant lost obligations.
States and transitions are enumeration units, not samples. The projection is a
stable-channel, one-intent model of four vertices; it is not a mechanised proof
of the full language, of arbitrary graphs or of implementation refinement
(\evpath{RESULTS.md}, section ``Finite projection'').

\paragraph{The binary-derived counterfactual.}
The counterfactual takes each episode's original obligation set and its
observed completion tags and asks what the source-window rule would have
discarded. It runs no modified binary and compares no competitor
(\evpath{RESULTS.md}, section ``Source-window counterfactual'').
\Cref{tab:eval:counterfactual} transcribes the strata of \evpath{summary.json}.

\begin{table}[ht]
\centering\footnotesize
\caption{Source-window counterfactual by stratum. Every figure is the
\code{episodes}, \code{source\_mutant\_losing\_episodes} or
\code{source\_mutant\_lost\_obligations} field of the corresponding
\code{strata} entry in
\protect\evpath{experiments/saner-2027/campaign-mechanisms-20260910/summary.json}; the
same figures appear in \protect\evpath{RESULTS.md} beside it. Units are episodes and
obligations.}
\label{tab:eval:counterfactual}
\begin{tabular}{@{}llrrr@{}}
\toprule
Stratum & Value & Episodes & Losing episodes & Lost obligations \\
\midrule
Class & fault matrix & 120 & 84 & 111 \\
      & no-fault control & 10 & 0 & 0 \\
      & foreign identity & 10 & 10 & 10 \\
\midrule
Boundary & build & 40 & 28 & 41 \\
         & before publication & 40 & 28 & 35 \\
         & after artefact write & 40 & 28 & 35 \\
         & foreign & 10 & 10 & 10 \\
         & none & 10 & 0 & 0 \\
\midrule
Channel & stable & 70 & 47 & 61 \\
        & beta & 70 & 47 & 60 \\
\midrule
Graph & chain & 28 & 20 & 38 \\
      & diamond & 28 & 20 & 25 \\
      & fan-out & 28 & 20 & 20 \\
      & two disconnected pairs & 28 & 14 & 14 \\
      & chain with shortcuts & 28 & 20 & 24 \\
\midrule
All & & 140 & 94 & 121 \\
\bottomrule
\end{tabular}
\end{table}

The figures the conference account quotes are the fault-matrix row: 111
obligations lost in 84 of the 120 fault episodes, with 28 of 40 episodes at
each of the three boundaries and 47 of 70 in each channel. The ten no-fault
controls lose nothing, which is the expected null. The ten foreign-identity
cases lose one obligation each and are reported separately because they test a
different intervention (\cref{sec:eval:rq2}). The counts are the sensitivity
of a deficient analytical rule under designed cases, not a win rate.

\emph{Answer to \evRQ{1}.} In every state the projection reached and in every
episode the binary completed, the consumer-relative window retained exactly the
residual set, and a source-relative window would have lost obligations in 260
of the 772 states and in 84 of the 120 fault episodes.

\subsection{RQ2: retry and identity stability}
\label{sec:eval:rq2}

\paragraph{The mechanism matrix.}
The matrix crosses five four-package topologies (chain, diamond, fan-out, two
disconnected pairs, chain with shortcuts) with stable and beta baselines, four
fault targets and three injection boundaries, giving 120 fault episodes; ten
no-fault controls and ten foreign-identity cases bring it to 140. Each fault
episode executes the same fault twice, repairs it and runs a final no-op; the
order is fixed by seed 20260910; build and publish concurrency are one
(\evpath{PROTOCOL.md}). The run took 113.653827~s, from 2026-09-09
20:39:05.988761 to 20:40:59.642604 UTC, and passed 140 of 140 episodes with no
failed, invalid or excluded case (\evpath{summary.json}, \code{execution} and
\code{executor}). \Cref{tab:eval:mechanism} gives the five checks.

\begin{table}[ht]
\centering\footnotesize
\caption{Mechanism checks of the 140-episode matrix, from the
\code{mechanism\_checks} object of
\protect\evpath{experiments/saner-2027/campaign-mechanisms-20260910/summary.json}.
The 130 denominator is the 120 fault episodes plus the ten foreign-identity
cases; the ten no-fault controls do not enter it (\protect\evpath{RESULTS.md}).}
\label{tab:eval:mechanism}
\begin{tabular}{@{}L{9cm}r@{}}
\toprule
Check & Passed / denominator \\
\midrule
Repeated fault preserved the residual plan and tags & 130 / 130 \\
Successful repair left an empty plan & 130 / 130 \\
Exactly one durable artefact write per obligation & 140 / 140 \\
Foreign artefact identity conflict was refused & 10 / 10 \\
Final no-op left the plan empty and tags unchanged & 140 / 140 \\
\bottomrule
\end{tabular}
\end{table}

Per stratum, 40 of 40 episodes passed at each boundary, 70 of 70 in each
channel and 28 of 28 on each graph (\evpath{summary.json}, \code{strata}). The
after-write injector lets the supplied local publisher write the artefact
durably and then return a nonzero exit; it simulates neither \code{SIGKILL}
nor a network partition nor power loss. The foreign-identity cases test that
adapter's immutable version-and-content identity, not an npm or OCI registry
(\evpath{summary.json}, \code{interpretation}). An independently implemented
offline checker reconstructed every plan from the raw command output, checked
tags, publisher events and the four stored artefact identities per case, and
re-enumerated the finite model; it reported no discrepancy over 140 episodes
(\evpath{summary.json}, \code{offline\_consistency\_check}). A negative check
on a corrupted copy failed as intended (\evpath{RESULTS.md}, ``Evidence
integrity'').

\paragraph{Reference agreement on the companion campaigns.}
The 1.10.0 companion campaign planned 300 generated histories, seeds 90001 to
90300, against the restricted core oracle: 300 accepted matching rows, 0
semantic rejections, 0 invalid executions, 0 rows with differences
(\evpath{experiments/saner-2027/campaign-1.10.0-20260909/summary.json},
\code{core}). The corrected 1.9.0 campaign reports the same 300 of 300 on the
same seeds (\evpath{paper/saner/evidence/EXPERIMENTS-1.9.0-20260909.md}). The
oracle is the historical model of the conference paper's restricted profile;
300 of 300 on a fixed sample is agreement on that sample, not a universal
conformance result.

The recovery stage of each companion campaign drives 20 episodes over graphs
of 20 and 100 packages with failure probabilities .1 and .3 and seeds 100 to
104. In 1.10.0 all 20 converged with progress on every run and at most four
runs, distributed as nine episodes in two runs, five in three and six in four,
over 251 raw command records with no invariant error (\evpath{summary.json},
\code{recovery}). In the corrected 1.9.0 run the same 20 episodes converged in
at most four runs over 248 command records, distributed nine, six and five
(\evpath{EXPERIMENTS-1.9.0-20260909.md}). The two campaigns are not pooled.
\evpath{runs/run-v1.10.0-linux-arm64-20260909/recovery-comparison.json} locates
the difference in one episode, graph size 100, probability .3, seed 100: the
first attempt left 20 packages pending in 1.10.0 against 18 in 1.9.0, so the
later seeded deny sets diverged and the remaining counts ran 20, 2, 1, 0
against 18, 2, 0, which accounts for one extra release, tag and status
triplet (\evpath{EXPERIMENTS-1.10.0-20260909.md}). The traces do not attribute
the difference to the 1.10.0 feature; fixed seeds do not fix a concurrent
executor's schedule.

\paragraph{The 1.8.2 full-profile campaign.}
The 7 September 2026 campaign on the released 1.8.2 binary planned 3,000
generated instances in four shards: core seeds 1 to 1000 and 5001 to 5500, and
full-profile seeds over the same ranges
(\evpath{output/experiments-v182-summary.json}, \code{E1.shards}). Both core
shards agreed exactly, 1,000 of 1,000 and 500 of 500, with 0 invalid
executions. The full profile deliberately generates regressive prereleases:
of its 1,500 instances, 503 were refused with a validated \diag{E195}
diagnostic and 997 were accepted; of the accepted, 546 matched the historical
oracle; four refused dossiers also had no row difference, so 550 comparisons
were diagnostic-clean; the divergent package rows numbered 1,914, classified
by the campaign's heuristic classifier as PG1 862, PG2 375, PG2x 14 and PG3?
663 with none unclassified (\evpath{output/experiments-v182-summary.json},
\code{E1.totals} and \code{E1.divergence\_classes};
\evpath{experiments/results/v182-20260907/README.md}). Per shard the
full-profile figures were 343 refusals, 353 accepted-clean and 356
diagnostic-clean rows with classes 572, 240, 9 and 442 for seeds 1 to 1000,
and 160, 193, 194 and classes 290, 135, 5 and 221 for seeds 5001 to 5500
(\evpath{E1.shards}). The canonical manifest binds the four shards to extractor
SHA-256 \evhash{3d5eef3ef2bf87de7876ddf6d04e852028b72091dd2417ce4943c96e48011e52}
(\evpath{experiments/results/v182-20260907/e1/canonical-manifest.json}). A
refused diagnostic plan is not an executable release plan, the oracle
implements the historical CCME 1 model, and the PG3? label is heuristic, so
these counts are taxonomy evidence rather than CCME 2.0.0 conformance
(\evpath{output/experiments-v182-summary.json}, \code{limits}).

The first 1.8.2 attempt is retained at
\evpath{experiments/results/v182-20260907-attempt1}: its strict exit checker
classified the expected \diag{E195} refusals as invalid executions, and the
corrected campaign reran every seed from scratch rather than patching the
records (\evpath{experiments/results/v182-20260907/README.md}). Seed 580 of
the core profile, which produced an invalid fixture under 1.8.1, succeeded
under 1.8.2 in a changed fixture environment (tmpfs and a pre-run \code{git
fsck}); planner source is unchanged between the two releases, so the success
is not a product fix (\evpath{README.md}; \code{E1.normalized\_changes\_from\_v181}
lists that one core seed and 503 full-profile seeds).

A separate check in the same campaign ran \SPEC{} Appendix~B vectors 133 to
138 against the released 1.8.2 binary: 6 of 6 passed
(\evpath{output/experiments-v182-summary.json}, \code{CCME2}). Vector 135
exercises the observable downward eligibility result and is not an independent
implementation of the upward audit, so the proof remains the support for
\cref{thm:duality}.

\paragraph{The historical differential campaigns.}
The oracle of E1 (\evpath{experiments/e1/oracle.py}) implements the model of
this report as it stood against the CCME 1.0.0 snapshot; the generator
(\evpath{experiments/e1/gen.py}) produces layered instances and the realiser
(\evpath{experiments/shared/realize.py}) turns each into a real repository
(\evpath{experiments/README.md}). On the released 1.4.0 binary the core
profile agreed 1,000 of 1,000 on seeds 1 to 1000 and 500 of 500 on seeds 5001
to 5500. The full profile over seeds 1 to 1000 left 412 instances clean and
classified every divergent row as PG1 409, PG2 240, PG2x 58 or PG3 390, and
over seeds 5001 to 5500 left 222 clean with PG1 182, PG2 135, PG2x 23 and PG3
200, with zero unclassified (\evpath{experiments/README.md}, ``v1.4.0
campaign, remaining results''; \evpath{experiments/results/compare-v140-v160.md}).
The first shard's 1,097 divergent rows are the four predicted classes and
none outside them: the 240 PG2 prerelease-train discharges and the 58 PG2x
suppression-lapse and cascade variants are \cref{ob:pre,ob:g5} observed at
scale; the 409 PG1 channel-entry patches are the deliberate one-patch rule of
\spec{11.4} (\diag{W204}), which the report model does not transcribe; and the
390 PG3 bump-axis liveness rows arise because the oracle implements the
companion article's compact admission, which drops clause~\ref{a:live}. The
1.6.0, 1.7.0 and 1.7.2 campaigns reproduce the per-seed rows on all 3,000
instances with only the version string differing
(\evpath{experiments/results/compare-v140-v160.md},
\evpath{v170/compare-v160-v170.md}, \evpath{v172/compare-v170-v172.md}). These
records validate the CCME 1.0.0 model against those binaries; they do not
validate CCME 2.0.0 admission, whose window is the subject of \evRQ{1}.

\paragraph{Orphan recovery and mid-release pushes.}
E2 drives a six-package npm fixture against a Verdaccio registry behind a
failing proxy that answers 502 to named PUTs
(\evpath{experiments/README.md}). In the 1.8.2 campaign both repetitions
recorded the failed publish untagged, a same-version recovery and a final plan
releasing nothing (\evpath{output/experiments-v182-summary.json}, \code{E2}).
E5's three fresh 1.8.2 cells passed 8 checks for the orphan scenario, 12 for
a clean concurrent push and 14 for a conflicting concurrent push
(\code{E5}); the two push scenarios distinguish a completed release tag from a
concurrent source change on selected interleavings, not on a systematic sweep.
Its nine competitor cells were not rerun and remain historical (\code{limits}).

\paragraph{Historical competitor observations.}
The competitor phases of 30 August 2026 pin the newest versions published on
that date: Lerna 10.0.1, Nx 23.1.2, Turborepo 2.10.12 with Changesets 3.0.1,
and release-please 17.11.2 (\evpath{experiments/README.md}, ``Competitor tool
pins''). Each phase was run twice with identical protocol-relevant transcript
lines. Lerna tagged at version time, before publish; the denied \code{cli}
left \code{cli@1.0.1} tagged and orphaned, and recovery was blocked by
\code{EUNCOMMIT} until manual cleanup. Nx also tagged before publish and
propagated a patch to all five dependents of a minor change; its recovery is
\code{nx release publish} again, which skips already-published versions by
registry query. Changesets tagged after publish, left no tag for the denied
package, recovered by registry query, and gated propagation by version range;
its release intent is a changeset file consumed at version time, so the plan
is not a function of Git history and cannot be replayed. Turborepo neither
versions nor publishes. release-please has no local publish step and was
observed only through the E3a replay (\evpath{experiments/README.md}, ``E2
extension results''). These are observations of the pinned configurations,
not properties of every mode those tools expose.

\paragraph{Historical replay.}
E3a replays five releases each of two public repositories under a
reconstructed layout and policy. The 1.8.2 campaign reproduced the established
profile: 36 of 42 released package-versions exact for
\code{google-cloud-node} and 7 of 64 for \code{conventional-changelog}
(\evpath{experiments/results/v182-20260907/README.md}, and \code{E3} of
\evpath{output/experiments-v182-summary.json}, which records
\code{metrics\_equal\_v181} true for all ten rows). The per-release rows are
in \evpath{experiments/results/e3a/gcn.json} (released 1, 22, 11, 7, 1 and
exact 1, 17, 11, 6, 1) and \evpath{experiments/results/e3a/cc.json} (released
20, 8, 9, 6, 21 and exact 2, 3, 1, 0, 1). Neither project used the modelled
intent language, and the mismatches may reflect the reconstruction as much as
the projects' conventions; the experiment shows the limit of policy transfer
and does not claim those releases followed CCME.

\emph{Answer to \evRQ{2}.} Across the designed graph, channel, target and
fault factors of the 140-episode matrix, recorded progress survived repeated
failure without version drift, repair emptied the plan, and the identity
check refused every foreign artefact. The restricted oracle agreed on all 300
histories of each companion campaign and on all 1,500 core instances of the
1.8.2 campaign; the 1,914 divergent full-profile rows and the 1,097 rows of the
historical campaign are retained and classified, and they mark where the
restricted models stop rather than where the implementation fails.

\subsection{RQ3: availability boundaries and CI}
\label{sec:eval:rq3}

The companion campaigns execute two registry topologies and one CI fixture.
The three-package topology is an npm package, a Docker base image and a Docker
application; the four-package extension is \code{assets} (npm), \code{ui}
(npm), \code{base} (Docker) and \code{app} (Docker), where \code{ui} reads the
predecessor's local \code{dist} output and performs no registry install
(\evpath{experiments/saner-2027/campaign-next-release-prepared/CAMPAIGN-PLAN.md}).
Each topology repeats five configurations six times in a frozen seeded order:
mixed local-output scheduling, fully serialised controls, a missing barrier,
an injected base-publication failure with repair, and an injected app-build
failure with repair.

\begin{table}[ht]
\centering\footnotesize
\caption{The 1.10.0 companion campaign. Executions, initial outcomes, repairs,
no-op reruns and predicate counts are the \code{mixed} and \code{integration}
objects of \protect\evpath{experiments/saner-2027/campaign-1.10.0-20260909/summary.json};
the intervals are the six-decimal figures of
\protect\evpath{paper/saner/evidence/EXPERIMENTS-1.10.0-20260909.md}, whose
nine-decimal originals are the \code{overlap} arrays of the same
\protect\evpath{summary.json}. Predicates are repeated verifier checks, not samples.}
\label{tab:eval:rq3}
\begin{tabular}{@{}L{3.1cm}rrrrrr@{}}
\toprule
Topology & Executions & \shortstack[r]{Initial\\successes} &
\shortstack[r]{Designed\\failures} & Repairs & \shortstack[r]{No-op\\reruns} &
Predicates \\
\midrule
Three-package & 30 & 12 & 18 & 12 & 24 & 522 \\
Four-package & 30 & 12 & 18 & 12 & 24 & 798 \\
CI fixture & 3 repetitions & & & & & 81 \\
\bottomrule
\end{tabular}
\end{table}

Per topology the 18 designed failures are six of each negative configuration;
all 12 repairable executions completed at unchanged planned versions; the six
missing-barrier cases intentionally have no repair; the other 24 completed
cases performed no build or publication on their final unchanged run
(\evpath{EXPERIMENTS-1.10.0-20260909.md}). In the three-package mixed fixtures
the base build overlapped npm publication in all six observations, by
0.110716 to 0.151483~s. In the four-package mixed fixtures the \code{ui} build
overlapped the deliberately delayed \code{assets} publication by 0.070128 to
0.073725~s, and the base build started 6.292152 to 6.393487~s before the
\code{ui} publication ended and completed before that publication began in
all six cases. The latter is an admissible schedule: a build that reads local
output need not wait for its provider's publication, and the fixture's delays
expose availability rather than measure throughput. The three CI repetitions
each cover baseline publication, a specification change, a consumer change,
an injected failing test, a missing required test, repair and an unchanged
run; the three failing-test and three missing-test negatives blocked the
release stages, the failing-test attempts wrote no release tag, repairs
propagated the specification release to its consumer, and unchanged runs
executed the selected tests without release stages.

\paragraph{The corrected 1.9.0 campaign.}
The completed 1.9.0 run, \evpath{run-v1.9.0-linux-arm64-20260909-corrected},
has the same 12, 18, 12 and 24 counts per topology, the same 522, 798 and 81
predicate counts, overlaps of 0.101969 to 0.128675~s in the three-package
fixtures and 0.068203 to 0.077008~s between \code{ui} and \code{assets}, with
the base build starting 6.298757 to 6.311191~s before \code{ui} publication
ended (\evpath{paper/saner/evidence/EXPERIMENTS-1.9.0-20260909.md}). Its
provenance records all 14 stages passed and a \code{COMPLETE} marker written
at 2026-09-09T10:48:01 UTC
(\evpath{runs/run-v1.9.0-linux-arm64-20260909-corrected/provenance.json}).

\paragraph{The preserved failed oracle attempt.}
The first 1.9.0 attempt, \evpath{runs/run-v1.9.0-linux-arm64-20260909}, ran
from 05:35:36 to 07:37:14 UTC, passed its first thirteen stages and failed
the fourteenth, \code{mixed-4node-verification}, with exit code 1
(\evpath{provenance.json} of that run, \code{status: FAILED}). Its verifier
log ends in the assertion ``\code{mixed-2: base local build overlaps ui
publication}'' (\evpath{mixed-4node-verification.log}): the original oracle
required the base-build and \code{ui}-publication intervals to overlap
strictly, and the traces showed the valid schedule in which the base build
completed before \code{ui} publication began. The predicate was corrected
before the full rerun to the dependency-relevant condition, that the base
build starts before \code{ui} publication ends, with overlap and
completed-before-start both recorded descriptively
(\evpath{CAMPAIGN-PLAN.md}). The failed attempt has no \code{COMPLETE} marker
and is not counted as passing evidence; its recovery stage had reached a
maximum of five runs over 251 raw command records where the corrected run
reached four over 248, which is why the two are not pooled
(\evpath{EXPERIMENTS-1.9.0-20260909.md}). The 1.10.0 campaign reuses the
corrected protocol byte for byte, protocol manifest SHA-256
\evhash{a3366e2fbdc590be7e8f6268bc7390146669cf7a0395837c2531a0af5554dcaa}, with
no further oracle change (\evpath{campaign-1.10.0-20260909/summary.json}).

A separately frozen diagnostics addendum in the 1.10.0 campaign passed 16 of
16 purpose-selected cases for the standalone diagnostics command; it is
engineering evidence outside the 14-stage campaign and its sample counts
(\evpath{EXPERIMENTS-1.10.0-20260909.md}).

\emph{Answer to \evRQ{3}.} In these fixtures, on one Linux ARM64 environment,
a configured barrier distinguishes local build availability from required
publication, and both a failed and an absent required test prevent release;
the campaign measures neither concurrency at scale nor onboarding effort.

\subsection{RQ4: planning cost}
\label{sec:eval:rq4}

The E4 timing cells plan generated workspaces of $P$ packages with $U$ intent
units, in three graph shapes, with an optional padding history of commits that
touch no package. \Cref{tab:eval:rq4} transcribes the eight cells of the 1.8.2
campaign.

\begin{table}[ht]
\centering\footnotesize
\caption{Planning cost, 1.8.2 campaign: the eight rows of
\protect\evpath{experiments/results/v182-20260907/e4/timing.csv}, verbatim. Seconds
are the median, minimum and maximum of the cell's repetitions; MiB is the
maximum resident set size. The container was capped at two CPUs and 4~GiB on a
contended Apple-silicon Docker host
(\protect\evpath{experiments/results/v182-20260907/README.md}).}
\label{tab:eval:rq4}
\begin{tabular}{@{}rlrrrrrrr@{}}
\toprule
$P$ & Shape & $U$ & Padding & Edges & Median s & Min s & Max s & MiB \\
\midrule
100 & layered & 200 & 0 & 185 & 0.017 & 0.017 & 0.017 & 15 \\
1,000 & layered & 200 & 0 & 2,002 & 0.250 & 0.235 & 0.252 & 21 \\
10,000 & layered & 200 & 0 & 19,919 & 19.272 & 18.992 & 20.000 & 92 \\
1,000 & chain & 200 & 0 & 999 & 0.254 & 0.241 & 0.258 & 26 \\
10,000 & chain & 200 & 0 & 9,999 & 19.354 & 18.933 & 19.645 & 118 \\
1,000 & star & 200 & 0 & 999 & 0.244 & 0.237 & 0.247 & 19 \\
1,000 & layered & 1,000 & 0 & 2,002 & 0.290 & 0.280 & 0.302 & 26 \\
1,000 & layered & 200 & 5,000 & 2,002 & 0.664 & 0.640 & 0.665 & 30 \\
\bottomrule
\end{tabular}
\end{table}

The medians rise from 0.017~s at 100 layered packages to 19.272~s at 10,000,
and the 10,000-package chain records 19.354~s and 118~MiB, the largest of
both columns. Multiplying the intent five-fold at 1,000 packages moves the
median from 0.250 to 0.290~s; adding 5,000 padding commits moves it to
0.664~s, the history term of \cref{sec:complexity} showing through. The
fixture generator uses \code{git commit-tree} for long padding histories only
after a small fixture matched the ordinary generator's commit IDs, refs,
messages and log output, and every fixture passed \code{git fsck}
(\evpath{experiments/README.md}, gotcha 18 and the paragraph following it).

\paragraph{The loaded-host repetitions.}
The earlier campaigns measured the same headline cells under a different
harness and report much larger figures: in the 1.4.0 campaign the headline
cells were re-measured in five containers each, giving about 0.6~s at 100
packages, 2.8 to 6.7~s at 1,000 and 110 to 149~s at 10,000, with the spread
attributed to bimodal container placement on the host, and the
1.6.0 comparison records medians of medians of 0.61, 6.49 and 143.1~s against
0.60, 6.43 and 142.4~s (\evpath{experiments/README.md}, ``v1.4.0 campaign,
remaining results''; \evpath{experiments/results/compare-v140-v160.md}). Those
figures are not comparable with \cref{tab:eval:rq4}: the runner then retained
every generated workspace until the container exited and divided the resident
set size twice, and it was rewritten before the 1.8.x campaigns
(\evpath{experiments/README.md}, gotcha 17). The two sets are reported side
by side as what each harness measured; neither is a cross-version benchmark.

The E4 recovery cells of the 1.8.2 campaign converged in 24 of 24 episodes
within four runs, 31 total runs in the first repetition and 32 in the second
(\evpath{output/experiments-v182-summary.json}, \code{E4.recovery}); the 1.4.0
campaign's 24 episodes converged with a maximum of five runs and the 1.6.0
campaign's with four (\evpath{compare-v140-v160.md}).

\emph{Answer to \evRQ{4}.} On the capped host the observed cost grows sharply
between the 1,000- and 10,000-package cells and with history length; the
design isolates neither a general asymptotic bound nor an end-to-end figure,
and the loaded-host repetitions show how far the same cells move with the
harness.

\subsection{RQ5: distributed execution}
\label{sec:eval:rq5}

The execution profile of \spec{28} places build and publish stages on worker
nodes that share a Git repository with the orchestrator as their only
transport. Two build graphs were measured on 22 September 2026, both on Google
Compute Engine~\cite{gce}: the Linux kernel image-build farm of
kerneltests.org~\cite{kerneltests}, whose cells are independent, and the LLVM
monorepo, whose cells consume one another's outputs. The specification's own
record of the implementation and of both measurements is the entry
``2026-09-22: Distributed execution implemented'' of
\evpath{/Users/yohimik/Projects/dispat/specs/ccme-spec/DESIGN-HISTORY.md},
which this subsection follows in substance and cites for the figures it
states. The claim the measurements support is narrow and is stated once here:
on these fixtures, once each, a distributed release took no longer than the
same recipe on one node. No speedup of the profile or of the engine is claimed
beyond that. The raw observations of both graphs are copied, unedited, into
\evpath{experiments/distributed-20260922/evidence/} (the evidence tree below),
whose \evpath{MANIFEST.json} records every file's origin and SHA-256 and whose
\evpath{CROSS-CHECK.json} traces the headline figures of the two fork reports
and the specification entry to those raw files: 39 checks, 37 agreeing, 0
disagreeing, and two additional kernel trials that are disclosed there and
reported below rather than counted as headline figures.

\paragraph{Acceptance oracles.}
Each fork carries its own comparison script, and a run counts only if the
script accepts it. The kernel fork's \evpath{dispat/compare.py} accepts a
measured run only if its cases, kernel revision, compiler and assembler
identities match the retained public reference exactly, every cell succeeded,
no configuration is missing or duplicated, and the elapsed time is no greater
than the reference (\evpath{dispat/README.md} and \evpath{dispat/PARITY.md} of
\evpath{/Users/yohimik/Projects/linux-build-test}). The LLVM fork's
\evpath{.ci/dispat/gce/compare.py} compares the install tree of every package
file by file between the single-node and the distributed checkouts, with the
install prefix masked, the executable bit and symlink targets recorded, and
the toolchain archive compared by member list, sizes and modes; it prints the
receipts' own fingerprints beside the masked comparison and exits nonzero on
any difference. Neither script measures speed on its own; the kernel gate
folds elapsed time in, the LLVM pair reports it beside the product comparison.

\subsubsection{The kernel farm}
\label{sec:eval:rq5:kernel}

\paragraph{References.}
The public reference for each architecture group is the build step of the
upstream builder on kerneltests.org for kernel v7.3-rc4 at
\evhash{93f51579e7df248780214094418f205253383cc5}, read from the public build
pages and retained in \evpath{dispat/gce/parity-20260922-reference.json} of
the fork and in
\evpath{experiments/distributed-20260922/evidence/kernel-gce/round1-codex/baseline.json}:
alpha 623~s (build 332),
hexagon 111~s (build 330) and m68k\_nommu 128~s (build 331), with
\code{alpha-linux-gcc (GCC) 16.2.0} and \code{m68k-linux-gcc (GCC) 16.2.0}
over GNU assembler 2.47.20260726 from the kernel.org crosstool
builds~\cite{crosstool}, and \code{Ubuntu clang version 18.1.3 (1ubuntu1)} for
hexagon. The exact case lines are four for alpha, three for hexagon and five
for m68k\_nommu. These are individual build-step durations on the farm's
own, different, hardware, not the wall time of a farm run. The farm's
configuration gives image builds an exclusive lock, one per machine
(\evpath{master/master.cfg}, unchanged in the fork from upstream
\code{336111e3}), and its maintainer reports that the build directory is a RAM
disk that holds a single build (correspondence of 23 September 2026). The
fork's first report describes the reference steps as shared with another
build; the configuration does not support that. The references bound
acceptance and are not compared numerically beyond the gate. For context only, the first fork report
records that the farm's 57 upstream builders took 23.0 machine-hours in
3~h~52~min of wall time for v7.3-rc4 on six machines
(\evpath{dispat/gce/REPORT.md}); that figure was read from public pages and
enters no comparison here.

\paragraph{The first run.}
The first real run, on the \code{dispat-prototype} branch with a worker-nodes
build at commit \code{1146c210}, three VMs (an \code{e2-standard-4}
orchestrator and two \code{n2-custom-14-28672} spot workers) and kernel.org
crosstool gcc 14.3.0, is recorded in \evpath{dispat/gce/REPORT.md}; its
orchestrator and worker logs are copied under
\evpath{experiments/distributed-20260922/evidence/kernel-gce/first-run}. It is
not parity evidence: the compiler is not the reference compiler, the exported
pass, fail and skip counts were not retained, the wrapper accepted the
script's exit code as success although the farm script accumulates errors and
exits zero, and the raw logs were not found in the fork itself
(\evpath{dispat/PARITY.md}). What it established was the coordination profile
and two engine defects. Four groups finished; \code{m68k} was lost when
Compute Engine preempted its worker 18~minutes into the build; the run was
restarted twice because the worker's idle timer stopped a node five seconds
after a 38-minute build and stopped a fresh node after 34~s. Claim-to-result
was 2,250~s for alpha against the 623~s reference, 324~s for hexagon against
111 and 246~s for m68k\_nommu against 128, ratios of 3.5, 2.3 and 1.8 of the
script body to the reference step; the probe round trip was 4~s, assignment
to claim 37~s cold and 3~s warm, and a result reached the orchestrator in 5~s
(\evpath{dispat/gce/REPORT.md}). The copied logs bear these figures out:
\evpath{gce-timings.py} run over \evpath{first-run/gce-logs} reads the
worker's claim and completion stamps as 00:04:47 to 00:42:20 for alpha,
01:07:28 to 01:12:54 for hexagon, 01:07:27 to 01:11:33 for m68k\_nommu and
01:11:40 to 01:11:54 for um, within three seconds of the report's
claim-to-result figures at the logs' one-second resolution
(\evpath{CROSS-CHECK.json}, the first-run entries).

\paragraph{Round 1: every candidate failed the gate.}
Round 1 (\evpath{dispat/gce/PARITY-20260922.md}) installed the reference
toolchains, ran the development binary of SHA-256
\code{a7172a7f\ldots} named in \cref{tab:eval:boundaries} on an
\code{e2-standard-2} orchestrator, a \code{c3-standard-22} worker (11 cores,
22 threads) and a \code{t2d-standard-8} worker (8 cores, 8 threads), and
measured a grouped control and four candidates, each once. The gate is GNU
\code{time} elapsed of \code{dispat release} from invocation to exit,
including source transport, assignment and claim polling, worker checkout, the
complete build script, result transport and reporting; VM provisioning,
toolchain setup and one-time source preparation are outside it.

\begin{table}[ht]
\centering\footnotesize
\caption{Kernel round 1, from the results table of
\protect\evpath{/Users/yohimik/Projects/linux-build-test/dispat/gce/PARITY-20260922.md};
the release seconds are also the \code{phase-timing.json} of each candidate in
\protect\evpath{experiments/distributed-20260922/evidence/kernel-gce/round1-codex/curated-linux-parity-evidence.tar.gz}.
Script seconds are the complete build-script body per cell. One trial per
candidate. The reference is the kerneltests.org build step.}
\label{tab:eval:kernel1}
\begin{tabular}{@{}L{3.0cm}L{5.4cm}rrL{2.4cm}@{}}
\toprule
Candidate & Placement and script seconds & Release s & Reference s & Gate \\
\midrule
split hexagon & c3: defconfig 69; t2d: tinyconfig 54, allnoconfig 76 &
170.07 & 111 & fail, elapsed only \\
split m68k\_nommu & c3: m5272 35, m5407 34, m5249 35; t2d: m5307 65, m5475 49
& 159.03 & 128 & fail, elapsed only \\
batched hexagon & c3: allnoconfig+tinyconfig 65; t2d: defconfig 147 & 165.30
& 111 & fail, elapsed only \\
batched m68k\_nommu & c3: first two cases 66; t2d: last three 167 & 196.23 &
128 & fail, elapsed only \\
\bottomrule
\end{tabular}
\end{table}

\code{compare.py} accepted the cases, toolchain identities and revision of
every candidate and rejected only the elapsed time. The cold grouped control
ran the three groups sequentially on the C3 worker in 1,797.75~s of release,
of which the three script bodies were 1,608~s (hexagon 130, m68k\_nommu 154,
alpha 1,324); alpha's script alone exceeded its 623~s reference, so no
removal of dispatch overhead could have carried it. The first control
transported the full 1.4~GB repository and 1.9~GB working tree, its source
phase taking 45~s of which the recorded push was 44.392~s; an image-only
export reduced the transport to about 512~KB and later pushes recorded 0.031
to 0.050~s. Per cell, warm kernel fetches took 0.446 to 0.571~s and worktree
checkouts 7.026 to 8.327~s; assignment to claim took 1 to 4~s and completion
to coordinator receipt 1 to 5~s, with a 0~s scheduler queue gap. About 89\% of
the control's elapsed time was inside the build scripts
(\evpath{PARITY-20260922.md}). The failed setup attempt
\code{batched-hexagon-setup-failure} is retained apart from the timed
candidates; distcc was prepared and never measured. The round's raw evidence is
the curated archive named in that file, SHA-256
\evhash{2fcc02ec891df09044b45a8b9d783015545e984abafe7eff745edac21cf08255}, 145
files, copied to
\evpath{experiments/distributed-20260922/evidence/kernel-gce/round1-codex/curated-linux-parity-evidence.tar.gz};
the release elapsed figures of \cref{tab:eval:kernel1}, the control's
1,797.75~s, alpha's 1,324~s of script, the per-group 229, 173 and 1,340~s and
the 44.392~s push are the \evpath{*/phase-timing.json} members of that
archive, as \evpath{CROSS-CHECK.json} traces them.

\paragraph{Round 2: calibration.}
Round 2 (\evpath{dispat/gce/PARITY-20260922-round2.md}) first measured the
farm's own script bodies outside dispat, on prepared source trees, with
\evpath{dispat/gce/calibrate.sh}; the raw lines are
\evpath{dispat/gce/round2/calibration/cal-c3.txt} and \evpath{cal-t2d.txt}, and
the installed toolchains are confirmed in
\evpath{cal-c3-toolchains.txt}. On \code{c3-standard-22}, hexagon defconfig
alone at 33 jobs took 73.92~s; the three hexagon configurations started
together at 11 jobs each took 120.11, 81.04 and 63.67~s, the group finishing
after 120.12~s; the five m68k\_nommu configurations started together at 7 jobs
each took 128.16, 127.41, 126.65, 127.04 and 103.16~s, the group finishing
after 128.17~s; alpha defconfig, allnoconfig and tinyconfig alone at 33 jobs
took 68.48, 33.16 and 24.04~s (\evpath{cal-c3.txt}). On \code{t2d-standard-8},
hexagon allnoconfig alone took 87.51~s and the two-batch m68k\_nommu split took
236.18 and 295.49~s (\evpath{cal-t2d.txt}); round 1 had measured hexagon
defconfig at 147~s there. Three facts followed: a single m68k\_nommu build
leaves part of the C3 machine idle, so several may share it; the three
hexagon configurations crowd each other on one machine; and the T2D machine
cannot hold hexagon defconfig inside the reference.

\paragraph{Round 2: topologies and results.}
The budget for the longest cell is the reference less the release's fixed
overhead, measured at 3 to 8~s from release start to \code{worker ready}:
102~s for hexagon and 119~s for m68k\_nommu. The topologies, generated into
\evpath{dispat/gce/round2/par-c3/dispat.hexagon.yaml} and
\evpath{dispat.m68k_nommu.yaml}, pin cells to the machine that starts the
release with \code{runOnly: orchestrator} and leave the rest to the worker.
For hexagon, \code{execution.concurrency} is 2: defconfig at 24 jobs and
tinyconfig at 9 jobs on \code{par-c3}, allnoconfig at 12 jobs on the
\code{t2d-standard-8} worker \code{w1}. For m68k\_nommu it is 4: the four
costliest configurations at 8 jobs each on \code{par-c3}, \code{m5475evb} at
12 jobs on \code{w1}. Both files carry \code{runOnly: [worker, orchestrator]}
at the top so that a worker's cell is built on the worker and recorded by the
orchestrator. \Cref{tab:eval:kernel2} transcribes every timed release of
the round: the rejected first candidate, the no-op rerun that planned
nothing, the four accepted trials and the two single-node runs. The fork's
files are copied unchanged under
\evpath{experiments/distributed-20260922/evidence/kernel-gce/fork-codex-kernel-parity-c699bf3d/dispat/gce/round2/},
and the paths below are relative to that directory in both places.

\begin{table}[ht]
\centering\footnotesize
\caption{Kernel round 2, every timed release. Elapsed is the \code{elapsed=}
field of the GNU \code{time} record in \protect\evpath{par-c3/*.time} and
\protect\evpath{solo-c3/*.time} (also \protect\evpath{*.elapsed}); script and
lease seconds per cell and the \code{compare.py} verdict with its
\code{speedup} ratio (reference over elapsed) are from
\protect\evpath{verdicts.txt}; all under \protect\evpath{dispat/gce/round2} of
the fork and under
\protect\evpath{experiments/distributed-20260922/evidence/kernel-gce/fork-codex-kernel-parity-c699bf3d/dispat/gce/round2}.
The reference is the kerneltests.org build step. Each row is one release.}
\label{tab:eval:kernel2}
\begin{tabular}{@{}L{2.4cm}L{4.9cm}rrL{3.3cm}@{}}
\toprule
Run & Cells: script s (lease s) & Elapsed s & Ref.\ s & Verdict, ratio \\
\midrule
hexagon, first candidate & c3: tinyconfig 60 (12.0), defconfig 84 (12.0);
w1: allnoconfig 80 (12.0) & 111.47 & 111 & reject, elapsed only;
0.9957836189109177 \\
hexagon, rerun of the unchanged configuration & no cell built: every package
reported \code{unchanged} & 3.83 & 111 & not compared; planned nothing \\
hexagon, trial 1 & c3: tinyconfig 70 (0.0), defconfig 85 (0.0); w1: allnoconfig
82 (0.0) & 95.89 & 111 & accept; 1.1575763896130984 \\
hexagon, trial 2 & c3: tinyconfig 71 (0.0), defconfig 84 (0.0); w1: allnoconfig
81 (0.0) & 92.22 & 111 & accept; 1.2036434612882239 \\
m68k\_nommu, trial 1 & c3: m5249evb 95 (1.0), m5272c3 96 (1.0), m5407c3 96
(14.0), m5307c3 97 (15.0); w1: m5475evb 50 (0.0) & 115.61 & 128 & accept;
1.1071706599775106 \\
m68k\_nommu, trial 2 & c3: m5249evb 105 (0.0), m5407c3 105 (0.0), m5307c3 106
(0.0), m5272c3 107 (0.0); w1: m5475evb 51 (0.0) & 115.48 & 128 & accept;
1.1084170419120194 \\
hexagon, one node alone & solo: tinyconfig 91 (13.0), allnoconfig 105 (13.0),
defconfig 110 (12.0) & 122.62 & 111 & reject, elapsed only;
0.9052356874898059 \\
m68k\_nommu, one node alone & solo: m5475evb 56 (0.0), m5407c3 91 (0.0),
m5307c3 92 (0.0), m5272c3 73 (58.0), m5249evb 72 (66.0) & 138.01 & 128 &
reject, elapsed only; 0.9274690239837694 \\
\bottomrule
\end{tabular}
\end{table}

\code{compare.py} accepted the exact case lines, the compiler, the assembler
and the revision in all seven releases that built cells and rejected three of
them on elapsed time alone; every cell reported \code{passed}, three cases for
hexagon and five for m68k\_nommu, with no skips and no unrecognised lines
(\evpath{verdicts.txt}). The eighth release, \evpath{par-c3/hexagon-2}, was a
rerun of the unchanged hexagon configuration before the release tags were
reset: it planned nothing, logged \code{unchanged} for every package and
exited in 3.83~s (\evpath{par-c3/hexagon-2.elapsed} and
\evpath{hexagon-2.release.log.gz}); it is a release with nothing to do and is
listed so that no run of the round is omitted. The gate was therefore met
twice for hexagon after one rejected candidate, at 95.89 and 92.22~s against
111, and twice for m68k\_nommu with no rejected candidate, at 115.61 and
115.48~s against 128. The rejected first candidate reached 111.47~s against
the 111~s reference with a 12.0~s lease wait on each of its three cells; it
preceded the leased-tree change described below. Their margins, 13 to 17\% and 10\%, exceed
the 3.7~s spread between the two hexagon trials and say nothing about other
groups (\evpath{PARITY-20260922-round2.md}, ``Limits'').
\evpath{dispat/gce/round2/timing.txt}, produced by \evpath{dispat/timing.py}
over the coordinator and worker logs, reports the same overheads from the
worker's side across the five paired runs: the initial source push at
0.000~s with 0.009 to 0.024~s of Git work, assignment to claim in 0 to 3~s,
the per-cell farm checkout in 0 to 1~s, a finished cell reaching the
coordinator in 0 to 5~s, and the release's own Git subprocesses at 7.205 to
7.231~s in every run.

\paragraph{What the harness changes bought.}
Three changes separate round 1's split hexagon at 170.07~s from round 2's
trial 2 at 92.22~s on the same two machine shapes
(\evpath{PARITY-20260922-round2.md}, ``What each harness change bought'').
First, a leased source tree replaced the per-cell worktree: a node prepares
one pristine tree per build slot during setup and a cell leases one under an
advisory lock after checking its revision and cleanliness, which cost 0.0~s in
every warm run against round 1's 7.026 to 8.327~s per cell. The check is not
free on a tree nothing has read: the first candidate recorded 12.0~s of lease
on all three cells and m68k\_nommu trial 1 recorded 14.0 and 15.0~s on two
slots used for the first time (\evpath{verdicts.txt}); the startup script now
performs that read once per slot, a change made after these runs and credited
with none of the figures. Second, recording a worker cell on the orchestrator:
in the first candidate the worker's build finished at $+106$~s and the release
exited at 111.47~s because the publish stage went back to the worker for one
more poll, a four-second authorisation and a report; \code{runOnly: [worker,
orchestrator]} removed 5.5~s from the end of every release, and trials 1 and 2
exit within a second of their last build. Third, a generator that writes the
topology down (\code{--orchestrator}, \code{--concurrency},
\code{--cell-jobs}, \code{--batch}, \code{--publish-on-orchestrator}) so that
placement, slots, jobs and batching are regenerated rather than edited. None of
the three is an engine change; the first is the fork's, the second a
configuration the engine already offered, the third tooling.

\paragraph{The same-hardware pairs.}
To say what the second machine is worth, each group ran once more on a single
\code{c3-standard-22} with no worker and nothing else changed: the same cells,
per-cell jobs, build-slot setting and recipe, with \code{runOnly:
orchestrator} and no \code{workers} key
(\evpath{dispat/gce/round2/solo-c3/dispat.hexagon.solo.yaml} and
\evpath{dispat.m68k_nommu.solo.yaml}). The single node reports the same
processor as \code{par-c3}, Intel Xeon Platinum 8481C at 2.70~GHz, 11 cores
and two threads each, in zone us-central1-b because the shape had become
unavailable in us-central1-a. Both single-node runs built every case correctly
and were rejected on elapsed time alone: hexagon at 122.62~s and m68k\_nommu
at 138.01~s, ratios 0.9052356874898059 and 0.9274690239837694 of the reference
(\evpath{verdicts.txt}). Against the paired trials the second machine, a
slower one of another family, is what carries both groups under the gate.

Two features of the single-node runs are stated rather than hidden. First,
\code{execution.concurrency} did not bound a release that names no worker:
hexagon was configured with two build slots and started all three cells at
once, m68k\_nommu with four and started all five; such a release emits no
assignment, claim or probe, which also removes the 3 to 8~s of preflight the
paired runs pay, so the single-node figures are the best the machine can do
with the recipe, not a throttled version of it. Second, the lease and
freshness check dominated the m68k\_nommu run: five cells checked five trees
at the same moment on one disk, and two recorded 58 and 66~s of
\code{source\_lease\_seconds} with build scripts of only 73 and 72~s
(\evpath{solo-c3/results/build-m68k_nommu-m5272c3_defconfig.g5NdBk/result.json}
and \evpath{build-m68k_nommu-m5249evb_defconfig.LEf7Hy/result.json}, field
\code{timings.source\_lease\_seconds}, in the fork's
\evpath{dispat/gce/round2} and in the evidence tree's copy); the three hexagon
cells recorded 12.0
to 13.0~s (\evpath{verdicts.txt}). The 138.01~s figure is therefore pessimistic
by some part of that cost. The conclusion does not rest on the single-node
runs alone: the calibration measured the same shape building each whole group
outside dispat and without any lease, the hexagon trio in 120.12~s and the
five m68k\_nommu cells in 128.17~s of script body
(\evpath{cal-c3.txt}), both at or above their references before a second of
release overhead.

\paragraph{Why these layouts do not transfer to the farm.} The paired
layouts ran two hexagon or four m68k\_nommu image builds at once on the
\code{c3-standard-22}, the single-node runs three and five, and the
calibration the same stacking, each build in its own directory under a
\code{/tmp} on the boot disk (the image's default; no mount output was
retained). The farm runs one image build per machine in a RAM disk that holds
one build, so none of these layouts fits its machines as they are, and the
measured build directories were disk-backed where the farm's are in memory.
The round shows what a second machine bought on this hardware; it does not
show that the farm would finish sooner.

\paragraph{alpha, and why it was not run.}
The alpha group was not attempted. Its \code{allmodconfig} is one
configuration of one cell and cannot be split. The calibration put the other
three alpha cases at 68.48, 33.16 and 24.04~s on \code{c3-standard-22}, 125.68~s
together, and round 1 had measured all four as one script body on the same
shape at 1,324~s; the subtraction gives \code{allmodconfig} a lower bound of
about 1,198~s, lower because round 1's small cases shared a build directory
with it. With about 9~s of fixed overhead the budget for the cell is about
610~s, a factor of 1.96 over what the largest machine the project's 32-vCPU
quota allows can do. The fork estimates that a \code{c3-standard-88} would
clear the budget at 60\% scaling efficiency and a \code{c3-standard-44} would
not at the measured 0.78; both are estimates and neither was provisioned
(\evpath{PARITY-20260922-round2.md}, ``alpha: the capacity this project does
not have''). alpha is therefore a group the experiment could not measure, not
one it measured and failed.

\paragraph{What the engine cannot express, and what it logged.}
The engine has two placement classes and no per-worker pinning; it places a
ready package on any free node without regard to speed. The two classes
sufficed here only because each group needed exactly two machines and the
class pinned to the release's own machine was the fast one; a second
\code{t2d-standard-8} worker could receive hexagon defconfig, at 147~s there,
and lose the group. In the m68k\_nommu trials the worker finished at $+56$ and
$+66$~s and sat idle for 59 and 49~s while the orchestrator finished four
cells; the fork estimates a perfect split near 100~s of build against the
107~s measured (\evpath{PARITY-20260922-round2.md}, ``The placement dispat
cannot express''). The worker logged one \diag{E226} ``assignment rejected,
unreadable'' warning per release for the snapshot branch, five lines in all
and the only warnings a worker produced; a release refused to plan without a
pushable \code{origin} for its lock tag, so source and coordination refs now
share the mailbox repository; a package \code{env} may not carry a key with the
\code{DISPAT\_} prefix, so the node-describing variables moved to the node's
environment; a second run of an unchanged configuration planned nothing and
exited in 3.83~s (\evpath{par-c3/hexagon-2.elapsed}, the row of
\cref{tab:eval:kernel2}), so repeated trials begin by removing the release
tags; and the preflight probe,
a poll round trip whose back-off reaches 5~s while a worker is idle, is the
largest fixed overhead left at 3 to 8~s. The round's five VMs lived 10 to
22~minutes each and were deleted; the list-price estimate for the round is
about USD~1.48 (\evpath{PARITY-20260922-round2.md}, ``Cost'').

\subsubsection{The LLVM pair}
\label{sec:eval:rq5:llvm}

\paragraph{Fixture and machines.}
The question is whether a dispat release of the fork's own recipe, with its
build stages delegated to workers, produces the same products as the same
release on one machine and takes no longer. The fork is
\code{yohimik/llvm-project}, branch \code{dispat-gce}, LLVM \code{main} at
\code{a9f4d3ff} (24.0.0git); the report is \evpath{.ci/dispat/gce/REPORT.md}
at the branch tip \code{952c2de4}, and every fork figure below is from it
unless another file is named. The recipe \evpath{.ci/dispat/build.py} runs with
\code{LLVM\_FAST=1} (assertions, examples, docs and benchmarks off),
\code{LLVM\_SKIP\_CHECKS=1} (no \code{check-*} targets), the X86 target only,
\code{LLVM\_BUILD\_JOBS=14}, \code{LLVM\_LINK\_JOBS=2} and
\code{LLVM\_TOOLCHAIN\_PROVIDERS=llvm,lld,polly}; the same values appear in
\evpath{.ci/dispat/gce/dispat.gce.yaml} (two workers, mailbox the
orchestrator's own bare \code{origin.git}) and \evpath{dispat.local.yaml} (the
same recipe with no \code{execution} object). Baselines
\code{dispat/<pkg>/v24.0.0} are tagged at \code{a9f4d3ff}, so every plan is
24.0.0 to 24.1.0. The machines are \code{lb-orch}, an \code{e2-standard-4}
with 80~GB, and \code{lb-w1} and \code{lb-w2}, \code{n2-custom-14-49152} (7
cores, 14 threads) with 150~GB, zone us-central1-a, on-demand, all on Ubuntu
24.04.5 with g++ 13.3.0, cmake 3.28.3 and ninja 1.11.1. The recipe's identity
record embeds the compiler, cmake and ninja versions, the source revision and
the recipe's own hash into every receipt (\evpath{build.py},
\code{identity}), and a consumer refuses a provider whose receipt does not
match the tree it installed (\code{verify\_provider}).

\paragraph{Nine attempts.}
\Cref{tab:eval:llvm-attempts} lists every attempt. Six have a retained
orchestrator log and five of those a GNU \code{time} record, copied from the
measuring session's scratchpad into
\evpath{experiments/distributed-20260922/evidence/llvm-gce/logs/}; the mapping
of file to attempt is the \code{llvm\_attempt\_mapping} object of
\evpath{experiments/distributed-20260922/evidence/MANIFEST.json}, and the wall
times below were read from those \code{.time} files. Attempts 1, 3 and 6 have
no orchestrator log; attempts 1 and 6 are covered only by cumulative worker
logs, attempt 3 by nothing, its log having been overwritten on the node
(\evpath{MANIFEST.json}).

\begin{table}[!htp]
\centering\footnotesize
\caption{The nine LLVM attempts, from the attempts table of
\protect\evpath{.ci/dispat/gce/REPORT.md} on \code{origin/dispat-gce} of
\protect\evpath{/Users/yohimik/Projects/llvm-project}. Wall times and exit statuses in
the last column are the \code{Elapsed (wall clock) time} and \code{Exit
status} lines of the named \code{release-*.time} file under
\protect\evpath{experiments/distributed-20260922/evidence/llvm-gce/logs/},
and the published and failed counts are the \code{done} line of the matching
\code{release-*.log}; a wall time with exit 1 includes the failed bundle step;
``none'' means no orchestrator log was retained.}
\label{tab:eval:llvm-attempts}
\begin{tabular}{@{}rL{1.5cm}L{6.6cm}L{4.4cm}@{}}
\toprule
\# & When (UTC) & What happened & Retained log; wall; exit \\
\midrule
1 & 03:54 to 04:38 & \code{llvm} built and published through \code{w1};
\code{lld} and \code{polly} refused its transported tree because the recipe's
fingerprint hashed full file and directory modes, which are the creating
process's umask. Discarded; the fingerprint now records the executable bit and
link targets only. & none \\
2 & 04:40 to 05:07 & \code{llvm}, \code{lld}, \code{polly} built and published
through the workers in 27~min~21~s; the \code{toolchain} bundle failed at once
on an orchestrator with no C++ compiler for the identity record. Discarded; a
provisioning omission. & \code{release-dist2}; no \code{.time}; \code{done}:
3 published, 1 failed \\
3 & 05:20 to 05:21 & Baseline attempt; the clone on \code{w2} lacked the
\code{dispat/*/v24.0.0} tags, so the plan started from the initials and the
recipe refused the major. Discarded; the tags were pushed. & none \\
4 & 05:45 to 06:12 & Distributed: \code{llvm} on \code{w1}, \code{lld} on
\code{w2}, \code{polly} on \code{w1}, all published; the bundle refused
\code{polly}'s transported tree while \code{llvm} and \code{lld} passed. &
\code{release-dist3}; 26:47.69; exit 1; 3 published, 1 failed \\
5 & 06:13 to 06:49 & Baseline on \code{w2} alone, four packages published:
\code{llvm} 18:03, \code{lld} and \code{polly} in parallel to 06:33:24, the
bundle 15:53 of single-threaded xz. & \code{release-local}; 36:09.09; exit 0;
4 published \\
6 & 06:51 to 07:14 & Distributed again with the recipe at \code{2b03be6a}
(directory records left out of the fingerprint). Interrupted when the fixture
changed to the matrix. & none \\
7 & 07:16 to 07:45 & The matrix, distributed: \code{llvm} on \code{w1} 22:23,
\code{llvm-aarch64} on \code{w2} 23:06, \code{lld} and \code{polly} to
07:43:28, four published; the bundle refused \code{polly} once more. Diagnosed
as a dispat defect: captured outputs passed through the checkout's
\code{.gitattributes} text conversion, which rewrote CRLF pairs inside three
libraries. & \code{release-dist5}; 28:45.63; exit 1; 4 published, 1 failed \\
8 & 07:46 to 08:27 & The matrix on \code{w2} alone: five published, the two
heavy builds sharing the node for 38:38 and 38:51, \code{lld} and \code{polly}
to 08:27:14 and 08:27:28, the bundle 18~s. & \code{release-local3};
41:12.02; exit 0; 5 published \\
9 & 08:29 on & The matrix, distributed, with the fixed dispat
\code{1.11.0-rc.4+main.272772cb}: five published. & \code{release-dist6};
27:48.94; exit 0; 5 published \\
\bottomrule
\end{tabular}
\end{table}

The single-heavy fixture of attempts 1 to 6 was replaced by the matrix because
one package held nine tenths of the work: the matrix adds
\code{llvm-aarch64}, the same sources built for the AArch64 target as a package
that depends on nothing (\evpath{.ci/dispat/packages/llvm-aarch64.yaml}), so
that a single node must serialise two heavy builds and a pool need not, and
it sets \code{LLVM\_BUNDLE\_COMPRESSION=none} on both sides because a
single-threaded xz of the whole install tree was a 16-minute serial step on
either side and says nothing about placement (\evpath{REPORT.md}, ``The matrix
fixture'').

\paragraph{The matrix result.}
\Cref{tab:eval:llvm-matrix} is the pair. The distributed release took 27:48.94
against 41:12.02 on one node, five packages published and exit 0 on both
sides, the same recipe and the same per-package settings. The distributed
side had been measured once before with the same placement, attempt 7, at
28:45.63 to the same point plus a bundle that failed on the dispat defect;
that attempt is not a second repetition of the accepted pair.

\begin{table}[!htp]
\centering\footnotesize
\caption{The LLVM matrix pair, attempts 8 and 9. Per-package figures are the
results table of \protect\evpath{.ci/dispat/gce/REPORT.md}; the whole-release wall
times are the \code{Elapsed (wall clock) time} lines of
\protect\evpath{release-local3.time} and \protect\evpath{release-dist6.time}
under \protect\evpath{experiments/distributed-20260922/evidence/llvm-gce/logs/},
and the \code{done} lines of the matching logs record 5 published and 0
failed on both sides. One repetition per side.}
\label{tab:eval:llvm-matrix}
\begin{tabular}{@{}L{4.6cm}L{3.6cm}L{5cm}@{}}
\toprule
& Baseline, \code{w2} alone & Distributed, \code{w1} + \code{w2} +
\code{lb-orch} \\
\midrule
\code{llvm} (X86) & 38:38, sharing the node & 22:08 on \code{w1}, including
the 1.4~GB result push \\
\code{llvm-aarch64} & 38:51, sharing the node & 23:22 on \code{w2}, including
its result push \\
\code{lld} + \code{polly} & 2:16 & 2:13 on the two workers, after \code{llvm}'s
outputs reached them \\
bundle (uncompressed tar) & 0:18 & 0:32 on \code{lb-orch} \\
whole release, GNU \code{time} & 41:12.02, exit 0 & 27:48.94, exit 0 \\
\bottomrule
\end{tabular}
\end{table}

The design-history entry adds the transfer volumes: \code{llvm}'s declared
outputs, 1.4~GB in 2,647 files, were pushed once and fetched by two consumers
and the orchestrator, \code{llvm-aarch64}'s about the same, \code{lld}'s
420~MB and \code{polly}'s 22~MB likewise, and the bundle's 1.85~GB never
travelled; \code{lld} took 1~min~47~s and \code{polly} 2~min~13~s on the
opposite workers once \code{llvm}'s outputs had reached them
(\evpath{DESIGN-HISTORY.md}, entry of 2026-09-22). What the pair shows is the
mechanism \spec{28.7} names as the profile's required gain, at work on a real
dependent graph: \code{lld} and \code{polly} consumed \code{llvm}'s outputs on
nodes that never built \code{llvm}.

\paragraph{Products.}
The bundle on the orchestrator passed the recipe's own integrity check for
every provider it merged: the receipt each worker wrote describes, file by
file, the tree the orchestrator installed, so the transport is byte-faithful.
\code{compare.py} then finds the two sides differ in bytes, and every
difference is of the reproducible-builds kind: every tool of equal size
differs first at its \code{.note.gnu.build-id} and then only in the compiled-in
VCS revision, since LLVM embeds the checkout's commit (\code{6c1a5fcd} on the
single node, the snapshot commit \code{dbc1e8a7} on a worker, which builds
from the prepared input state rather than from the trigger commit) and the
build id is a hash over that content; \code{llvm-config} and \code{lld} also
embed the repository string, which differs in length between the two
checkouts; \code{LLVMPolly.so}, \code{libPolly.a}, \code{libPollyISL.a} and
\code{liblldCommon.a} differ by a few kilobytes of RPATH and \code{\_\_FILE\_\_}
strings that carry the checkout's absolute path; and \code{libLLVMCore.a},
which embeds none of these, is byte-identical (\evpath{REPORT.md},
``Products''). Byte identity throughout is a recipe matter,
\code{-DLLVM\_APPEND\_VC\_REV=OFF}, a file prefix map and a relative RPATH,
which this recipe does not set. The claim the pair supports is therefore: the
same five products, built by the same recipe on two placements, verified by
the recipe's own receipts, differing only in the revision, build-id and path
strings the recipe chooses to embed.

\paragraph{The single-heavy negative case.}
Attempts 4 and 5 are the case \spec{28.7} predicts against distribution. On
\code{w2} alone, \code{llvm} built in 18:03, \code{lld} and \code{polly} in
2:13 in parallel, and the three were published at 20:16; the xz bundle then
took 15:53 for a whole release of 36:09.09. Distributed, \code{llvm} took
22:11 on \code{w1} including its 1.4~GB result push, \code{lld} and
\code{polly} 2:07 and 2:19 on \code{w2} and \code{w1}, and the three were
published at 26:47.69, a wall time that includes the bundle step failing on
the orchestrator's missing compiler and ends in exit 1 (\evpath{REPORT.md},
``The single-heavy fixture''; \evpath{release-local.time} and
\evpath{release-dist3.time} in the evidence tree). With one
package holding nine tenths of the work, moving a 1.4~GB install tree to
consumers that build in two minutes cannot pay for itself; the matrix is the
fixture that says what the profile is for.

\paragraph{What the pair measured about the engine.}
A cold worker's first claim comes minutes after a warm one's, because the
node's poll fetches the snapshot branch before it can read the assignment,
about 10~minutes for the 3.6~GB history on a cold cache and 6~s warm.
Coordination is cheap beside the builds: probe 1 to 4~s, claim 3 to 6~s warm,
a result noticed within 5~s, authorisation to publish command 8 to 12~s.
Placement is first-free, and the pair's workers are equal, so nothing here
says what an unequal pool would do; the kernel rounds do
(\evpath{REPORT.md}, ``What the pair measured about dispat''). The
design-history entry records the earlier LLVM coordination measurements the
implementation was tuned on: a probe answering in about 4~s, an assignment
claimed 3~s after it is written on a warm node and 37~s on a cold one, a
result noticed within 5~s, and a first snapshot transport that carried the
repository's whole history, 7.3~million objects, in 17 to 20~minutes from a
small orchestrator and was paid on every run until the repository's own
remote became the mailbox (\evpath{DESIGN-HISTORY.md}, entry of 2026-09-22).

\paragraph{Three defects found and fixed.}
The pair found one defect in the fork's recipe and two in the engine. The
recipe's integrity fingerprint hashed full file and directory modes, which are
the creating process's umask; a worker built directories 0775 and the engine
installs files as 0644 or 0755 and recreates directories, so a faithfully
transported tree failed its receipt (attempt 1). The fingerprint now records
the executable bit and symlink targets only, and directories not at all
(\evpath{build.py}, \code{fingerprint\_records}); attempt 4's diagnosis
confirmed that every one of the 124 entries \code{w1} captured was byte- and
mode-identical to the file the orchestrator installed. In the engine, a result
carrying outputs reported under a fixed 30~s bound that killed the 1.4~GB
push; the fix, at \code{f0c9e959} on the feature branch, ``give a result that
carries build outputs the transfer window to report'', makes the report use
\code{execution.transfer.timeout} (commit subject from \code{git log} of
\evpath{/Users/yohimik/Projects/dispat}; \evpath{REPORT.md}). And captured
outputs passed through the checkout's \code{.gitattributes}: \code{polly}'s
tree is marked as text, and \code{git add} rewrote every CRLF pair inside
three binaries, so the bundle refused the tree in attempts 4 and 7; the fix on
\code{main} at \code{9f649f88}, ``capture build outputs as the build wrote
them, whatever the checkout's attributes say'', hashes captured outputs with
no filter and fills the index by hand, and \code{272772cb} adds the
integration test that reads a library with CRLF pairs through a checkout
marked as text (\code{git log}; \evpath{DESIGN-HISTORY.md}).

\paragraph{The binaries.}
Attempt 9 ran \code{1.11.0-rc.4+main.272772cb}, the version string its
release log prints (\evpath{release-dist6.log}, field \code{dispat}). The
\code{binaries} object of \evpath{MANIFEST.json} records the deployed binary
by its SHA-256 prefix \code{feba0210d13c90fa} only, with the note that it was
built in a checkout carrying uncommitted documentation edits, so that Go
stamped it \code{vcs.modified=true}, and that it differs from a clean stamped
build only by that stamp; it is not reproducible from the commit alone. The
reproduction kit's recipe
\evpath{experiments/distributed-20260922/docker/dispat/Dockerfile} rebuilds
the pinned binaries from clean checkouts, and the hashes it wrote are in
\evpath{docker/dispat/verified-20260922/*/dispat.sha256}: \code{f0c9e959} unstamped gives
\evhash{63441cc66f68d07b5efbf9a42edd5559304d153d48b82b1dfb037349e20606ed},
byte-identical to the binary of kernel round 2 and LLVM attempts 1 to 8;
\code{272772cb} stamped gives
\evhash{46bc2cf821dd57479c27bf2a4b7eee63b69e01db6f18fbd6e663ad66e6812ba8}
and unstamped
\evhash{f1dec45d48df4d7fe5a681f5244ce946d71e140ccfbe97c2f8f3a6842f69b196}.
The version string in the log and the clean rebuild together are the pin this
report can offer for attempt 9. The reference engine attributed a unit's whole source set to every reached dependent until product \code{main} \code{e461d7d9}, after the measured build \code{272772cb}; the fixtures measured in \cref{sec:eval:rq5} carry single-source units only, which both rules attribute identically. The defect and its fix are \cref{ob:attribution}.

\paragraph{Caveats, as the specification states them.}
The design-history entry qualifies both pairs, and the qualifications are
repeated here without softening. One repetition per side on both graphs, so
variability is unmeasured; the kernel's second node is a slower machine of
another family, so each kernel pair is one node against that node plus a
second, slower one; the kernel groups have only independent cells to
distribute and do not exercise the transfer of dependent outputs; the LLVM
matrix's two heavy builds are one project for two targets, a fair matrix but
not two projects; the LLVM distributed side had two build nodes against one;
and the single-node m68k\_nommu figure includes 58 to 66~s of source-lease
contention on two of its cells (\evpath{DESIGN-HISTORY.md}, entry of
2026-09-22). Two details of the entry disagree with the retained files and
are recorded as observed on 22 September, the product side having been told:
the entry dates the LLVM pair 2026-09-23, while every retained log and
\code{.start} file is stamped 2026-09-22 UTC, 07:16 to 08:58, which no time
zone of the author's moves to the next day; and it names the fixture commit
\code{6c1a5fcd}, where the fork's own commit of the matrix fixture is
\code{d8fb8408} and the report's is \code{952c2de4}. The driver record
\evpath{llvm-gce/drivers/pair4.out} in the evidence tree explains the second:
the recipe change was applied on the orchestrator as a patch
(\evpath{llvm-gce/drivers/0001-feat-llvm-aarch64-build-the-AArch64-target-as-a-pack.patch})
and the resulting tip there, \code{6c1a5fcd}, is the commit under test that
the workers built, so the two hashes name the same change on two checkouts,
as the fork report says of its earlier pair (\code{fbc4dcf8} on the
orchestrator, \code{022f7284} locally, ``same tree''). Tree identity for the
matrix pair was not verified here, because \code{6c1a5fcd} exists only on the
deleted machine.

\paragraph{Public reference timings, as context only.}
Two public figures were gathered on 21 September 2026 and are retained; they
were read from public pages, concern other hardware and other recipes, and are
not compared numerically with the pairs. The upstream
\code{llvm/llvm-project} workflow run 34549173201 of 2026-09-11 ran 294.9166666666667
minutes from start to last update, and its \code{ubuntu-22.04} release-package
job took 154.1 minutes on a runner labelled \code{depot-ubuntu-22.04-16}, of
which 95.51666666666667 minutes were the build step and 56.35 the tests
(\evpath{output/saner-2027/llvm-forecast-20260921/timing-summary.json},
\code{runs[0]}; the README beside it records the same job as 154.10 minutes
and notes that it corrects an earlier approximate figure of 153). That run
builds the full release binaries with tests; the pair builds a 40-minute
fixture with no checks, so the two are not the same work. For the kernel farm,
the figure of 23.0 machine-hours in 3~h~52~min on six machines is already
cited above from \evpath{dispat/gce/REPORT.md}, with the same caveat.

\emph{Answer to \evRQ{5}.} On the kernel fixture the distributed release met
the public reference gate twice for each of two groups (hexagon 95.89 and
92.22~s against 111, after one rejected candidate at 111.47~s; m68k\_nommu
115.61 and 115.48~s against 128, with no rejected candidate) where the
same recipe on the same single node did not (122.62 and 138.01~s), each
single-node run once; alpha was outside the project's capacity and was not
run. On the LLVM fixture the distributed release of a five-package matrix took
27:48.94 against 41:12.02 on one node, produced the same five products under
the recipe's receipts, and the negative case predicted by \spec{28.7} was
observed at 26:47.69 against 20:16. On these fixtures, once each, distributed
took no longer than single node; nothing more is claimed.

\subsection{Threats to validity}
\label{sec:eval:threats}

\paragraph{Mechanism matrix and finite projection.}
The matrix is a designed, deterministic set of 140 episodes, not a sample of
production releases; the counts are exhaustive over the design and carry no
population inference. The after-write injector is an adapter response, not a
crash; the identity check is the supplied local publisher's, not a registry's.
The projection is four vertices, one intent and the stable channel; the
counterfactual runs no modified binary. The runner image cannot be rebuilt
bit-for-bit on a fresh machine because \code{apk add} consults an external
repository, so an exact rerun needs the locally present image
\code{sha256:6f92932f\ldots} (\evpath{campaign-mechanisms-20260910/README.md}).

\paragraph{Companion campaigns.}
The oracle of the 300-history stage is the restricted core profile and cannot
support a refinement claim; the recovery, topology and CI stages report
repeated verifier predicates, which are not independent observations; the
artificial delays in the mixed fixtures expose availability and are not
throughput; the environment is one Linux ARM64 Docker host and validates
nothing about native Windows. The failed 1.9.0 attempt shows that a verifier
can encode a wrong expectation, which is why its correction is documented
rather than the attempt deleted.

\paragraph{1.8.2 and historical campaigns.}
The E1 oracle is the CCME 1.0.0 model, so its divergence classes are taxonomy
against that model and the PG3? label is heuristic; a refused diagnostic plan
is not a release plan. The E3a replay reconstructs layouts and policies the
projects never declared. The competitor observations are of pinned versions in
one configuration each. Planning-cost figures are host-sensitive, taken on a
capped and contended host, and the earlier harness measured the same cells
very differently; none is a cross-version or cross-tool benchmark. The seed
580 change between 1.8.1 and 1.8.2 is not attributable to the product.

\paragraph{Kernel rounds.}
One pair of machine shapes, one region, one kernel revision, two trials per
accepted topology and one single-node run per group; the gate compares
against public build-step times measured on the farm's own, different,
hardware, where each image build has its machine to itself and builds in a RAM
disk, while the measured layouts stacked two to five builds on one machine in
disk-backed directories. The single-node
runs are penalised by the cold-metadata lease cost on two m68k\_nommu cells and
freed of the preflight the paired runs pay; the calibration outside dispat is
the corroboration. alpha, every QEMU group and the 57-group matrix are
unmeasured; distcc is unmeasured; the raw evidence of round 1 is a curated
archive copied into the evidence tree rather than files in the fork, and the
first run's logs, which the fork's \evpath{PARITY.md} could not find, are in
the evidence tree's \evpath{first-run} directory.

\paragraph{LLVM pair.}
One repetition per side; the distributed side had two build nodes against
one; the two heavy builds are the same sources for two targets; the recipe
omits the check targets and the assertions and is not the fork's release
recipe; six of nine attempts failed or were discarded for the reasons the table gives
and the accepted pair is attempts 8 and 9 only; the deployed binary carried
\code{vcs.modified=true} and is recorded by hash prefix only; and the raw
logs were copied from the measuring session's scratchpad into the evidence
tree, where \evpath{MANIFEST.json} records their origin, and are not in the
fork.

\subsection{Reproduction}
\label{sec:eval:repro}

The Docker and Terraform kit for the distributed measurements is at
\evpath{experiments/distributed-20260922/}: \evpath{README.md} and
\evpath{pins.json}; the \evpath{docker/} recipes (\evpath{dispat},
\evpath{kernel-node}, \evpath{llvm-node}); \evpath{terraform/kernel} and
\evpath{terraform/llvm} with a shared node module; \evpath{scripts/kernel/run.sh}
and \evpath{scripts/llvm/run.sh}, which drive a bare fleet to
\evpath{verdicts.txt} and \evpath{compare.txt} respectively; and a
\evpath{local/} Compose check. The \evpath{evidence/} directory beside them
holds the copied raw observations with their \evpath{MANIFEST.json} and
\evpath{CROSS-CHECK.json}.

What the kit pins, and what was verified on 22 September 2026 on an Apple
silicon workstation with Docker Desktop 29.6.1 and Terraform 1.16.3. The
experiment binaries rebuild byte for byte from their commits in a
digest-pinned \code{golang:1.26.5-alpine} with the flags of the original
builds: \code{f0c9e959} unstamped gives
\evhash{63441cc66f68d07b5efbf9a42edd5559304d153d48b82b1dfb037349e20606ed},
the binary the kernel round-2 report names; \code{272772cb} stamped
\code{1.11.0-rc.4+main.272772cb} gives
\evhash{46bc2cf821dd57479c27bf2a4b7eee63b69e01db6f18fbd6e663ad66e6812ba8},
the clean form of the LLVM matrix binary, whose deployed copy differs by Go's
\code{vcs.modified} stamp alone (records under
\evpath{docker/dispat/verified-20260922/}). The node images take Ubuntu 24.04
by digest, every package from the Ubuntu archive snapshot
\code{20260922T000000Z}, and the kernel.org crosstool 16.2.0 tarballs by
checksum; the kernel node image reports the identities the fork's oracle
requires (Ubuntu clang 18.1.3, GCC 16.2.0 with GNU assembler 2.47.20260726).
The four round-2 configurations are regenerated by the fork's own generator
with flags that reproduce the retained files byte for byte apart from the
mailbox path. Both Terraform roots validate and plan against the original
projects (thirteen resources each: network, subnet, two firewall rules, two
keys, the signing secret, three nodes and their secret files); the machine
shapes, zones, disk sizes, image and lifetimes are those of the measured
fleets. \code{terraform apply} was not run by the author of this section, so
the cloud pairs were not re-measured here; the kit is the recipe, and its
README states the expected verdicts and the cost.

The pipeline itself was run end to end on the workstation: two containers of
the kernel node image, an orchestrator and a worker sharing a mailbox
directory, the exported harness at the fork's commit, the generated
\code{hexagon} configuration, one \code{dispat release}. The worker built
\code{allnoconfig}, the orchestrator \code{defconfig} and \code{tinyconfig},
the release exited 0 and the fork's \code{compare.py} accepted the run: three
cases, the kernel revision \code{93f51579}, the compiler identity, and 63.47 s
of elapsed time (\evpath{experiments/distributed-20260922/runs/local-kernel-20260922T184954Z/verdict.json}).
That elapsed time is a laptop's on an arm64 toolchain and says nothing about
the parity of \cref{sec:eval:rq5:kernel}; what the run shows is that the
placement, the mailbox transport, the cell harness and the oracle work as
wired, from the pinned images alone.
